% answers_short.tex —— 1995 年·上海卷·文史类 · 简短答案速查
% 由 tools/make_short_answers.py 自 answers.tex 的 \gksolhead/\gkanswer 提取，
% 并与 1995/1995_数学真题_上海文_答案要点.md 逐题比对后生成，勿手改（改源 MD 后重跑）。

\begin{gknotice}
\textbf{使用说明}：本节只给最终答案便于快速核对，完整推导见解析版。答案为编写者独立推导，\textbf{非官方标准答案}。
\end{gknotice}

\gkpart{选择题}
\begin{center}\begin{tabular}{ccccccccc}
\toprule
题号 & 1 & 2 & 3 & 4 & 5 & 6 & 7 & 8 \\
\midrule
答案 & C & B & B & C & A & A & D & B \\
\bottomrule
\end{tabular}\end{center}

\gkpart{填空题}
\begin{itemize}[leftmargin=2.6em]
  \item[\textbf{9.}] $x<-3$ 或 $x>4$。
  \item[\textbf{10.}] $3x+4y=0$ 和 $3x-4y=0$。
  \item[\textbf{11.}] $y=54.8\left(1+\dfrac{x}{100}\right)^8$。
  \item[\textbf{12.}] $1$。
  \item[\textbf{13.}] $11$。
  \item[\textbf{14.}] $y^2=-8x+8$。
  \item[\textbf{15.}] $y=-\sqrt{\dfrac{x-1}{3}}$，$x\ge1$。
  \item[\textbf{16.}] $(\lg2,+\infty)$。
  \item[\textbf{17.}] $\left[-\dfrac{3\pi}{2},\dfrac\pi2\right]$。
  \item[\textbf{18.}] $A_1$组$-2<r<0$。$B_1$组答案：$e$。
  \item[\textbf{19.}] $A_1$组$350$ 种。$B_1$组答案：$\dfrac{25}{33}$（或 $0.76$）。
  \item[\textbf{20.}] $A_1$组$37.70\,\mathrm{dm}^3$；原评分标准接受 $37.68$ 或 $37.69$。$B_1$组答案：$14$ 天。
\end{itemize}

\gkpart{解答题}
\begin{itemize}[leftmargin=2.6em]
  \item[\textbf{21.}] $A_2$组$-\dfrac45$。$B_2$组答案：（1）$\overrightarrow{OD}=\left(0,-\dfrac12,\dfrac{\sqrt3}{2}\right)$；（2）$\cos\theta=-\dfrac{\sqrt{10}}5$。
  \item[\textbf{22.}] $A_2$组$(z_1,z_2)=\left(\dfrac{\sqrt3}{2}+\dfrac12i,-\dfrac{\sqrt3}{2}+\dfrac12i\right)$，或互换。$B_2$组答案：$f(x)=(x+1)^2$；所围面积为 $\dfrac13$。
  \item[\textbf{23.}] $V=2\sqrt3$；二面角大小为 $\arctan\dfrac{2\sqrt3}{3}$。
  \item[\textbf{24.}] （1）$S=\dfrac{4m^2n^2\tan\theta}{n^2+m^2\tan^2\theta}$；（2）$n<m$ 时 $u=2mn$，$m<n$ 时 $u=\dfrac{4m^2n^2}{m^2+n^2}$。题面还允许 $m=n$；此时两式均给 $u=2m^2$，分段边界在原答案中未列。
  \item[\textbf{25.}] （1）$f(x)=x^2-(t+2)x+t+1$；（2）$m=t^2+3t+3$，$n=-t^2-3t-2$；（3）$q=\dfrac{\sqrt{13}-2}{\sqrt{13}+2}$，$S=\pi\left(2+\dfrac{17\sqrt{13}}{26}\right)$。
\end{itemize}
